%==============================================================================
% CHAPTER 6 --- EVALUATION AND RESULTS
%==============================================================================
\chapter{Evaluation and Results}
\label{chap:evaluation}

This chapter defines the evaluation of the AntiGravity system across the four dimensions defined in Chapter~\ref{chap:methodology}: retrieval quality, generation quality, system performance, and user acceptance. Each section presents the result structure for the relevant hypotheses; the quantitative values are to be filled in upon completion of the experiments ([TBC] placeholders), and their significance --- including where results fall short of targets --- will be interpreted accordingly.

%------------------------------------------------------------------------------
\section{Evaluation Setup}
\label{sec:eval-setup}
%------------------------------------------------------------------------------

\subsection{Experimental Conditions}

Three conditions are compared throughout:

\begin{itemize}
\item \textbf{Condition A --- LLM Alone:} Mistral 7B Instruct, temperature=0.1, no retrieval. The model answers from parametric knowledge only.
\item \textbf{Condition B --- Naive RAG:} FAISS flat retrieval (top-$k=5$, no re-ranking, no MCP layer), same prompt template, same model.
\item \textbf{Condition C --- RAG+MCP (AntiGravity):} Full system --- MCP-mediated retrieval, ACL filtering, cross-encoder re-ranking, Mistral 7B generation.
\end{itemize}

All three conditions use identical generation parameters (temperature=0.1, top\_p=0.9, max\_new\_tokens=512, repetition\_penalty=1.1) and the same underlying model. Observed differences between conditions are therefore attributable to the architecture, not to model variation.

\subsection{Benchmark Dataset}

The evaluation benchmark comprises 50 question-answer pairs derived from the DNSI SharePoint corpus following the protocol of Section~\ref{sec:method-evaluation}. Table~\ref{tab:benchmark-composition} shows the distribution across query categories.

\begin{table}[H]
\centering
\caption{Benchmark dataset composition}
\label{tab:benchmark-composition}
\begin{tabularx}{\textwidth}{@{}lXrr@{}}
\toprule
\textbf{Category} & \textbf{Description} & \textbf{Count} & \textbf{\%} \\
\midrule
Factoid & Single-document, specific fact retrieval & 20 & 40\% \\
Multi-hop & Synthesis required across 2+ documents & 15 & 30\% \\
Procedural & Step-by-step instruction queries & 10 & 20\% \\
Unanswerable & Answer not in corpus & 5 & 10\% \\
\midrule
\textbf{Total} & & \textbf{50} & \textbf{100\%} \\
\bottomrule
\end{tabularx}
\end{table}

Two independent annotators label the ground-truth relevant chunks for each question. Overall inter-annotator agreement is measured using Cohen's $\kappa$ \citep{Cohen1960Kappa} [TBC]; pairs on which the annotators disagree are adjudicated and revised before inclusion in the benchmark.

% [NOTE FOR AUTHOR: Insert actual Cohen's kappa value here, e.g. kappa = [VALUE],
%  and specify how many pairs were revised.]

%------------------------------------------------------------------------------
\section{Retrieval Evaluation}
\label{sec:eval-retrieval}
%------------------------------------------------------------------------------

\subsection{Results}

Retrieval is evaluated on Conditions B and C (Condition A has no retrieval component). Results are reported over the answerable questions in the benchmark. A BM25 lexical baseline \citep{Robertson2009BM25} provides a reference point for assessing the contribution of dense semantic search.

\begin{table}[H]
\centering
\caption{Retrieval evaluation results (answerable questions)}
\label{tab:retrieval-results}
\begin{tabularx}{\textwidth}{@{}lXXXXX@{}}
\toprule
\textbf{System} & \textbf{P@1} & \textbf{P@3} & \textbf{P@5} & \textbf{MRR} & \textbf{nDCG@5} \\
\midrule
BM25 (lexical baseline) & [TBC] & [TBC] & [TBC] & [TBC] & [TBC] \\
Cond. B: Naive RAG (no rerank) & [TBC] & [TBC] & [TBC] & [TBC] & [TBC] \\
Cond. C: RAG+MCP (with rerank) & \textbf{[TBC]} & \textbf{[TBC]} & \textbf{[TBC]} & \textbf{[TBC]} & \textbf{[TBC]} \\
\midrule
$\Delta$ C vs. B (absolute) & [TBC] & [TBC] & [TBC] & [TBC] & [TBC] \\
$p$-value (Wilcoxon, $\alpha=0.05$) & [TBC] & [TBC] & [TBC] & [TBC] & [TBC] \\
\bottomrule
\end{tabularx}
\end{table}

% [NOTE FOR AUTHOR: Replace all [TBC] values with actual experimental results.
%  Run the retrieval evaluation pipeline on the 50-item benchmark dataset and record
%  P@1, P@3, P@5, MRR, nDCG@5 for BM25, Naive RAG, and RAG+MCP.
%  Run Wilcoxon signed-rank tests for pairwise comparisons.]

\paragraph{Hypothesis H3 (dense retrieval > BM25):} Dense retrieval is expected to outperform BM25 on this corpus, where technical concepts are expressed in both French and English and frequently through acronyms, creating a vocabulary mismatch that lexical matching cannot resolve. Results will be reported upon completion of the evaluation pipeline.

\paragraph{Hypothesis H6 (re-ranking improves MRR by at least 15\%):} Cross-encoder re-ranking is expected to deliver measurable MRR improvement over bi-encoder-only retrieval, particularly on multi-hop queries. Statistical significance will be assessed via the Wilcoxon signed-rank test.

\subsection{Retrieval by Query Type}

\begin{table}[H]
\centering
\caption{Retrieval Precision@5 disaggregated by query category}
\label{tab:retrieval-by-type}
\begin{tabularx}{\textwidth}{@{}lXXX@{}}
\toprule
\textbf{Query Type} & \textbf{Cond. B} & \textbf{Cond. C} & \textbf{$\Delta$} \\
\midrule
Factoid ($n=20$) & [TBC] & [TBC] & [TBC] \\
Multi-hop ($n=15$) & [TBC] & [TBC] & [TBC] \\
Procedural ($n=10$) & [TBC] & [TBC] & [TBC] \\
\bottomrule
\end{tabularx}
\end{table}

% [NOTE FOR AUTHOR: Replace [TBC] with measured values per category.]

The expected pattern is that re-ranking produces its largest gains on multi-hop queries, which require retrieving semantically distant chunks that the bi-encoder alone ranks below more lexically similar but less relevant passages. This hypothesis will be confirmed or disconfirmed by the disaggregated results.

\subsection{Retrieval Failure Analysis}

% [NOTE FOR AUTHOR: After running the evaluation, perform manual inspection of failed
%  retrievals (zero relevant results in top-5 for Condition C) and classify them by
%  failure mode. Expected failure categories based on system design:
%  (1) queries containing DNSI-internal acronyms absent from the embedding model vocabulary;
%  (2) queries targeting documents whose PDF extraction produced no parseable text (scanned pages);
%  (3) queries referencing documents added to SharePoint after the most recent ingestion cycle.
%  Report the count per category and the remediation pathway for each.]

%------------------------------------------------------------------------------
\section{Generation Quality Evaluation}
\label{sec:eval-generation}
%------------------------------------------------------------------------------

\subsection{RAGAS Metrics}

Table~\ref{tab:ragas-results} reports RAGAS \citep{Es2023RAGAS} generation quality across all three conditions on the answerable questions.

\begin{table}[H]
\centering
\caption{RAGAS generation quality results (answerable questions)}
\label{tab:ragas-results}
\begin{tabularx}{\textwidth}{@{}lXXXX@{}}
\toprule
\textbf{Condition} & \textbf{Faithfulness} & \textbf{Answer Relevancy} & \textbf{Context Precision} & \textbf{Context Recall} \\
\midrule
Cond. A --- LLM Alone & [TBC] & [TBC] & N/A & N/A \\
Cond. B --- Naive RAG & [TBC] & [TBC] & [TBC] & [TBC] \\
Cond. C --- RAG+MCP & \textbf{[TBC]} & \textbf{[TBC]} & \textbf{[TBC]} & \textbf{[TBC]} \\
\bottomrule
\end{tabularx}
\end{table}

% [NOTE FOR AUTHOR: Replace all [TBC] with measured RAGAS scores.
%  Run ragas.evaluate() on each condition's (question, answer, contexts, ground_truth) tuples.
%  Condition A has no retrieval context, so Context Precision and Recall do not apply.
%  Report 95% confidence intervals alongside point estimates.]

\paragraph{Hypothesis H1 (RAG+MCP reduces hallucination vs. LLM Alone):} The fundamental grounding mechanism of RAG --- injecting relevant, verified document context into the LLM prompt at query time --- is expected to substantially reduce the generation of facts not present in the DNSI corpus. This hypothesis will be assessed via the RAGAS faithfulness metric and the human-annotated hallucination rate.

\paragraph{Hypothesis H2 (RAG+MCP faithfulness $\geq 0.85$ on benchmark):} NFR10 specifies a faithfulness threshold of 0.85. Whether this threshold is met will be determined by the RAGAS evaluation.

\subsection{Hallucination Rate Analysis}

Hallucination rate is assessed through human annotation of a random sample of generated responses ($n_{\text{sample}}$ per condition), using the taxonomy of \citet{Ji2023}: intrinsic hallucinations (contradictions of retrieved context) and extrinsic hallucinations (claims absent from any retrieved chunk).

\begin{table}[H]
\centering
\caption{Human-annotated hallucination rate by condition}
\label{tab:hallucination-rate}
\begin{tabularx}{\textwidth}{@{}lXXX@{}}
\toprule
\textbf{Condition} & \textbf{Intrinsic Halluc.} & \textbf{Extrinsic Halluc.} & \textbf{Total Rate} \\
\midrule
Cond. A --- LLM Alone & [TBC] & [TBC] & [TBC] \\
Cond. B --- Naive RAG & [TBC] & [TBC] & [TBC] \\
Cond. C --- RAG+MCP & [TBC] & [TBC] & [TBC] \\
\bottomrule
\end{tabularx}
\end{table}

% [NOTE FOR AUTHOR: Fill in measured hallucination rates. Use two independent annotators
%  and report inter-annotator agreement (Cohen's kappa). Provide sample size n per condition.]

\subsection{Hallucination Taxonomy}

% [NOTE FOR AUTHOR: After annotation, classify observed hallucinations by type and discuss
%  the dominant patterns per condition. Expected findings based on system design:
%  Condition A hallucinations should be primarily extrinsic (fabrication of DNSI-specific
%  facts not present in training data). Conditions B and C should show fewer extrinsic
%  hallucinations but may exhibit intrinsic hallucinations when retrieved context is ambiguous.
%  Condition C is expected to show the lowest overall rate due to better context quality
%  from re-ranking and structured citation enforcement in the prompt template.]

\subsection{Qualitative Response Analysis}

% [NOTE FOR AUTHOR: Include two to three example query-response pairs that illustrate
%  key differences between conditions. Select examples that demonstrate:
%  (1) a case where Condition A hallucinated and Condition C did not;
%  (2) a case where Condition B retrieved a wrong passage and Condition C corrected it
%      via re-ranking;
%  (3) a case where none of the conditions could answer correctly (retrieval failure).
%  Anonymize any DNSI-specific content before inclusion per the confidentiality agreement.]

%------------------------------------------------------------------------------
\section{System Performance Evaluation}
\label{sec:eval-performance}
%------------------------------------------------------------------------------

\subsection{End-to-End Latency}

System latency is measured as the elapsed time from query submission to complete response display, under a representative workload (sequential queries, no concurrency). Each component's contribution is profiled separately to identify the dominant bottleneck.

\begin{table}[H]
\centering
\caption{End-to-end latency breakdown by component (P50 / P95)}
\label{tab:latency-breakdown}
\begin{tabularx}{\textwidth}{@{}lXXX@{}}
\toprule
\textbf{Component} & \textbf{P50 (ms)} & \textbf{P95 (ms)} & \textbf{NFR Budget} \\
\midrule
Query embedding (bi-encoder) & [TBC] & [TBC] & 100 ms \\
FAISS retrieval ($k=20$) & [TBC] & [TBC] & 500 ms \\
ACL resolution (permission lookup) & [TBC] & [TBC] & 500 ms \\
Cross-encoder re-ranking ($k=20 \to 5$) & [TBC] & [TBC] & 1000 ms \\
MCP protocol overhead & [TBC] & [TBC] & 200 ms \\
LLM generation (Mistral 7B, CPU) & [TBC] & [TBC] & --- \\
\textbf{Total end-to-end} & [TBC] & [TBC] & \textbf{5000 ms (NFR01)} \\
\bottomrule
\end{tabularx}
\end{table}

% [NOTE FOR AUTHOR: Instrument the pipeline with timing probes at each component boundary
%  and record latency distributions over a representative query set (recommended: 50 queries).
%  LLM generation on CPU is expected to dominate at 25--45 s per query with Q4-quantized
%  Mistral 7B, substantially exceeding NFR01 (5 s). This is a known infrastructure
%  constraint, not an architectural limitation. Document the measured GPU estimate
%  (e.g. from llama.cpp benchmark with GPU offloading) to show the NFR01 path.]

\paragraph{NFR01 compliance:} The CPU-only deployment of Mistral 7B is expected to produce generation latencies in the range of 25--45 seconds for typical query-length responses at approximately 3--8 tokens/second, substantially exceeding the 5-second NFR01 target. This constraint is infrastructural: the DNSI server used for prototype deployment lacks GPU resources. GPU deployment is the required intervention to achieve NFR01 compliance; the architectural components outside the LLM generation step (embedding, FAISS, ACL, re-ranking, MCP overhead) are each expected to remain well within their individual NFR budgets.

\subsection{MCP Overhead Isolation}

\paragraph{Hypothesis H4 (MCP latency overhead $<$ 200~ms P95):} The MCP protocol adds a communication layer between the conversational client and the context server. This overhead is the sum of serialization/deserialization cost and network round-trip time on the Docker bridge network. The 200 ms budget defined in the non-functional requirements represents the maximum acceptable overhead for this layer.

% [NOTE FOR AUTHOR: Measure MCP overhead by comparing retrieval latency in two configurations
%  that differ ONLY in the transport: (A3) dense retrieval + re-ranking via direct function
%  call vs. (A4) the same retrieval + re-ranking pipeline through the MCP transport.
%  The difference isolates MCP protocol overhead from search and re-ranking time.
%  Report P50 and P95 and state whether H4 is confirmed or refuted.]

\subsection{Memory Footprint}

% [NOTE FOR AUTHOR: Report peak RSS memory consumption for each Docker service at steady
%  state under load:
%  - mcp-server: FAISS index in RAM + chunk metadata store
%  - streamlit-app: Mistral 7B GGUF model loaded by llama-cpp-python
%  This is relevant for deployment planning and confirms feasibility on the DNSI server.]

%------------------------------------------------------------------------------
\section{User Evaluation}
\label{sec:eval-user}
%------------------------------------------------------------------------------

\subsection{Participant Profile}

% [NOTE FOR AUTHOR: Report participant demographics and profile:
%  - N (number of participants, from DNSI employee pool)
%  - Job functions represented (e.g., IT architect, project manager, business analyst)
%  - Prior experience with SharePoint search (years of use, frequency)
%  - Prior experience with AI assistants
%  These data come from the demographic questionnaire administered before the SUS/TAM session.]

\subsection{SUS Results}

The System Usability Scale (SUS) is administered to all participants after a single-session evaluation using the five task scenarios of Appendix~\ref{app:user-study}. SUS scores are computed using the standard formula of \citet{Brooke1996SUS} and interpreted against the adjective rating scale of \citet{Bangor2009SUSAdjective}.

\begin{table}[H]
\centering
\caption{SUS scores for AntiGravity}
\label{tab:sus-results}
\begin{tabularx}{\textwidth}{@{}lXX@{}}
\toprule
\textbf{Metric} & \textbf{Value} & \textbf{95\% CI} \\
\midrule
Mean SUS Score & [TBC] & [TBC] \\
Adjective Rating & [TBC] & --- \\
$N$ (participants) & [TBC] & --- \\
\bottomrule
\end{tabularx}
\end{table}

% [NOTE FOR AUTHOR: Replace [TBC] with measured values. Also report the distribution
%  across SUS items to identify specific usability problems. SUS scores between 68 and 80
%  are classified as "good"; below 68 as "ok"; above 80 as "excellent".]

\subsection{TAM Results}

TAM \citep{Davis1989TAM} Perceived Usefulness (PU) and Perceived Ease of Use (PEOU) are measured using the instruments in Appendix~\ref{app:user-study}. SharePoint search satisfaction is collected as a baseline using the same 7-point Likert scale applied to the PU items.

\begin{table}[H]
\centering
\caption{TAM scores: AntiGravity vs. SharePoint baseline}
\label{tab:tam-results}
\begin{tabularx}{\textwidth}{@{}lXXX@{}}
\toprule
\textbf{Measure} & \textbf{AntiGravity} & \textbf{SharePoint Baseline} & \textbf{$p$-value} \\
\midrule
Perceived Usefulness (PU, /7) & [TBC] & [TBC] & [TBC] \\
Perceived Ease of Use (PEOU, /7) & [TBC] & --- & --- \\
\bottomrule
\end{tabularx}
\end{table}

% [NOTE FOR AUTHOR: Use a one-sample or paired Wilcoxon test to compare AntiGravity PU
%  to SharePoint baseline PU. Report effect size (Cohen's d or r) alongside p-value.]

\paragraph{Hypothesis H5 (user acceptance SUS $\geq$ 68 and TAM PU significantly higher than SharePoint baseline):} A SUS score of 68 corresponds to the empirical average (50th percentile) across published SUS studies \citep{Sauro2011SUS}, and scores approaching 71 map to the ``Good'' adjective rating of \citet{Bangor2009SUSAdjective}; SUS $\geq$ 68 therefore represents the minimum acceptable adoption threshold. The TAM comparison tests whether users perceive AntiGravity as more useful than the current SharePoint search workflow.

\subsection{Task-Based Evaluation}

Five representative task scenarios (Appendix~\ref{app:user-study}) are administered to each participant under both systems (counterbalanced order). Task completion and time-on-task are recorded.

\begin{table}[H]
\centering
\caption{Task-based evaluation results}
\label{tab:task-eval}
\begin{tabularx}{\textwidth}{@{}lXXX@{}}
\toprule
\textbf{Task} & \textbf{AntiGravity Completion} & \textbf{SharePoint Completion} & \textbf{Mean Time Saving} \\
\midrule
T1 (Factoid) & [TBC] & [TBC] & [TBC] \\
T2 (Multi-hop) & [TBC] & [TBC] & [TBC] \\
T3 (Procedural) & [TBC] & [TBC] & [TBC] \\
T4 (Multi-hop) & [TBC] & [TBC] & [TBC] \\
T5 (Unanswerable) & [TBC] & [TBC] & N/A \\
\midrule
\textbf{Overall} & [TBC] & [TBC] & [TBC] \\
\bottomrule
\end{tabularx}
\end{table}

% [NOTE FOR AUTHOR: Record task completion (binary: completed within 5-minute time limit)
%  and time-on-task (seconds). For T5, correct completion is the user correctly identifying
%  that the system cannot answer the question. Report time saving only for successful tasks.]

\subsection{Qualitative Feedback}

% [NOTE FOR AUTHOR: Synthesize themes from the semi-structured interviews (Appendix~\ref{app:user-study}).
%  At minimum, identify: (1) most valued features (expected: source citations, natural language query);
%  (2) primary adoption barriers (expected: response latency on CPU deployment);
%  (3) trust factors identified by users. Present as thematic summary with representative quotes,
%  anonymized per the DNSI confidentiality agreement.]

%------------------------------------------------------------------------------
\section{Ablation Study}
\label{sec:eval-ablation}
%------------------------------------------------------------------------------

The ablation study decomposes the contribution of each architectural component by evaluating four sub-configurations:

\begin{itemize}
\item \textbf{A1 --- LLM Alone} (= Condition A): no retrieval.
\item \textbf{A2 --- Dense retrieval, no re-ranking, no MCP}: bi-encoder + FAISS flat retrieval, direct function call (no MCP transport), top-$k=5$ injected into prompt.
\item \textbf{A3 --- Dense retrieval + re-ranking, no MCP}: same as A2 with cross-encoder re-ranking.
\item \textbf{A4 --- Full AntiGravity (= Condition C)}: A3 with MCP transport and ACL filtering.
\end{itemize}

The progression A1 $\to$ A2 isolates the contribution of dense retrieval. A2 $\to$ A3 isolates re-ranking. A3 $\to$ A4 isolates the MCP layer.

\begin{table}[H]
\centering
\caption{Ablation study results: P@5 and RAGAS Faithfulness}
\label{tab:ablation}
\begin{tabularx}{\textwidth}{@{}lXXXX@{}}
\toprule
\textbf{Configuration} & \textbf{P@5} & \textbf{Faithfulness} & \textbf{$\Delta$ P@5 vs. prev.} & \textbf{$\Delta$ Faith. vs. prev.} \\
\midrule
A1 --- LLM Alone & N/A & [TBC] & --- & --- \\
A2 --- Dense only & [TBC] & [TBC] & [TBC] & [TBC] \\
A3 --- Dense + rerank & [TBC] & [TBC] & [TBC] & [TBC] \\
A4 --- Full system & [TBC] & [TBC] & [TBC] & [TBC] \\
\bottomrule
\end{tabularx}
\end{table}

% [NOTE FOR AUTHOR: Replace all [TBC] with measured values.
%  The expected finding for A3 vs A4 (MCP contribution) is that P@5 and Faithfulness
%  differences are small or non-significant --- MCP's primary contributions are security
%  and standardization, not retrieval quality. This is an expected and valid finding;
%  interpret it as such in the discussion rather than treating it as a null result.]

%------------------------------------------------------------------------------
\section{Hypothesis Testing Summary}
\label{sec:eval-hypotheses}
%------------------------------------------------------------------------------

\begin{table}[H]
\centering
\caption{Summary of hypothesis tests and outcomes}
\label{tab:hypothesis-summary}
\begin{tabularx}{\textwidth}{@{}llXX@{}}
\toprule
\textbf{ID} & \textbf{Hypothesis} & \textbf{Metric} & \textbf{Outcome} \\
\midrule
H1 & RAG+MCP reduces hallucination vs. LLM Alone & Hallucination rate & [TBC] \\
H2 & RAG+MCP faithfulness $\geq$ 0.85 & RAGAS Faithfulness (C) & [TBC] \\
H3 & Dense retrieval $>$ BM25 & P@5 & [TBC] \\
H4 & MCP overhead $<$ 200~ms P95 & MCP latency P95 & [TBC] \\
H5 & SUS $\geq$ 68 and TAM PU $>$ SharePoint & SUS, TAM & [TBC] \\
H6 & Re-ranking improves MRR by $\geq$15\% & MRR (C vs. B) & [TBC] \\
\bottomrule
\end{tabularx}
\end{table}

% [NOTE FOR AUTHOR: Replace [TBC] with Confirmed / Partially Confirmed / Refuted
%  and add the key measured value in parentheses, e.g. Confirmed (P@5=0.XX, p<0.05).]

%------------------------------------------------------------------------------
\section{Chapter Summary}
\label{sec:eval-summary}
%------------------------------------------------------------------------------

% [NOTE FOR AUTHOR: Once all experiments are complete, write a 2--3 paragraph summary
%  that synthesizes the key findings across the four evaluation dimensions:
%  retrieval, generation, performance, and user acceptance. Highlight which hypotheses
%  were confirmed or refuted, and identify the primary gap between the designed system
%  and the production-ready target (expected: LLM inference latency on CPU deployment).
%  This summary feeds directly into the Discussion chapter (Chapter~\ref{chap:discussion}).]
